\section{Accelerated tail bounds and the remaining end-to-end lemma}
\label{sec:accelerated-work}

\subsection{Objective-gap handoff}

The objective-gap switch is chosen so that the nonaccelerated coordinate tail
fits inside the accelerated target budget.

\begin{theorem}[Accelerated LocSOR tail from an objective-gap handoff]
\label{thm:objective-handoff-tail}
Let \(x_J\) satisfy
\begin{equation}
    f(x_J)-f(x^\star)
    \leq
    \Delta_{\rm sw}
    =\frac{\alpha^{3/2}\eps}{1+\alpha}.
    \label{eq:objective-handoff-condition}
\end{equation}
Run the local SOR tail with \(\omega=1\).  Then
\begin{equation}
    \Work_{\rm tail}
    \leq\frac{1}{\sqrt\alpha\,\eps}.
    \label{eq:objective-handoff-work}
\end{equation}
\end{theorem}

\begin{proof}
Set \(\omega=1\) in \cref{eq:general-tail-work} and substitute
\cref{eq:objective-handoff-condition}:
\[
    \Work_{\rm tail}
    \leq
    \frac{(1+\alpha)[f(x_J)-f(x^\star)]}
         {\alpha^2\eps^2}
    \leq
    \frac{1}{\sqrt\alpha\eps}.
\]
\end{proof}

\begin{proposition}[AESP stage count to the handoff]
\label{prop:aesp-handoff-stage}
Let \(U_t\) be the AESP upper bound in
\cref{eq:aesp-objective-bound}, and let \(J\) be defined by
\cref{eq:switch-stage}.  Write
\begin{equation}
    \log_+(z):=\max\{0,\log z\}.
    \label{eq:log-plus}
\end{equation}
Then a sufficient stage count is
\begin{equation}
    J+1
    \leq
    1+
    \left\lceil
    \frac{1}{\rho_A}
    \log_+\!\left(
       \frac{200(1-\alpha)(1+\alpha)^2}
            {\alpha^{5/2}\eps}
    \right)
    \right\rceil,
    \label{eq:handoff-stage-count}
\end{equation}
where \(\rho_A=0.9\sqrt{\alpha/(1-\alpha)}\).  In particular,
\begin{equation}
    J=\mathcal O\!\left(
       \frac{1}{\sqrt\alpha}
       \log\frac{1}{\alpha\eps}
    \right).
    \label{eq:handoff-stage-asymptotic}
\end{equation}
\end{proposition}

\begin{proof}
The condition \(U_J\leq\Delta_{\rm sw}\) is
\[
    (1-\rho_A)^{J+1}
    \leq
    \frac{\alpha^{5/2}\eps}
         {200(1-\alpha)(1+\alpha)^2}.
\]
Use \((1-\rho_A)^k\leq e^{-\rho_A k}\), solve for \(k=J+1\),
and use the definition of \(\rho_A\).
\end{proof}

The handoff target is substantially coarser than the final PPR tolerance.
Equivalently, if one invokes the standard AESP theorem with a coarse solution
parameter \(\eps_{\rm burn}\) satisfying
\begin{equation}
    \frac{\alpha\eps_{\rm burn}^2}{2}
    =\frac{\alpha^{3/2}\eps}{1+\alpha},
    \qquad
    \eps_{\rm burn}
    =\sqrt{\frac{2\sqrt\alpha\eps}{1+\alpha}},
    \label{eq:coarse-aesp-tolerance}
\end{equation}
then its objective-gap target is exactly the switch scale.  This is the formal
sense in which the hybrid uses only AESP's early or coarse-accuracy regime.

\subsection{Trajectory-dependent full complexity}

Let
\begin{equation}
    B_J:=\sum_{t=1}^{J}\Work_t^{\rm in},
    \qquad
    \Lambda_J:=\max_{1\leq t\leq J}\Lambda_t.
    \label{eq:burn-work-and-locality}
\end{equation}
The truncated version of the AESP work analysis gives
\begin{equation}
    B_J
    =\softO\!\left(\frac{\Lambda_J}{\sqrt\alpha}\right),
    \label{eq:truncated-aesp-work}
\end{equation}
where the soft notation hides the outer geometric schedule and the logarithmic
inner contractions in \cref{eq:aesp-stage-work}.

\begin{theorem}[Trajectory-dependent hybrid work]
\label{thm:trajectory-total}
The proof-oriented hybrid satisfies
\begin{equation}
    \Work_{\rm hybrid}
    \leq
    \softO\!\left(\frac{\Lambda_J}{\sqrt\alpha}\right)
    +\frac{1}{\sqrt\alpha\eps}.
    \label{eq:trajectory-total-work}
\end{equation}
\end{theorem}

\begin{proof}
Add the truncated AESP work \cref{eq:truncated-aesp-work} to the tail bound
\cref{eq:objective-handoff-work}.
\end{proof}

\begin{theorem}[Conditional accelerated local bound]
\label{thm:conditional-total}
Suppose the early AESP stages satisfy
\begin{equation}
    \Lambda_t\leq\frac{K}{\eps}
    \qquad\text{for every }1\leq t\leq J,
    \label{eq:early-locality-assumption}
\end{equation}
where \(K\) is independent of graph size.  Then
\begin{equation}
    \Work_{\rm hybrid}
    =\softO\!\left(
       \frac{K+1}{\sqrt\alpha\eps}
    \right).
    \label{eq:conditional-total-work}
\end{equation}
In particular, if \(K=\mathcal O(1)\), the desired scale
\(\softO(1/(\sqrt\alpha\eps))\) follows.
\end{theorem}

\begin{proof}
Insert \cref{eq:early-locality-assumption} into
\cref{eq:trajectory-total-work}.
\end{proof}

\begin{remark}[The exact missing lemma]
The existing AESP theorem controls outer objective convergence and bounds each
stage by its realized locality factor.  It does not prove
\cref{eq:early-locality-assumption} on every graph.  Thus the full theorem is
not yet unconditional.  The scientific core of the project is to prove a
useful version of this early-locality statement, weaken it to a summable
boundary overhead, or exhibit a counterexample that identifies the necessary
algorithmic safeguard.
\end{remark}

\subsection{A stronger residual-mass handoff}
\label{sec:mass-handoff}

The weighted original-gradient mass
\begin{equation}
    M(x):=\norm{D^{1/2}g(x)}_1
    \label{eq:mass-functional}
\end{equation}
controls both active volume and the work of Gauss--Seidel cleanup.

\begin{lemma}[Mass controls the final active-set volume]
\label{lem:mass-volume}
For
\(
S_\eps(x):=\{u:\abs{g_u(x)}\geq\alpha\eps\sqrt{d_u}\}
\),
\begin{equation}
    \vol(S_\eps(x))
    \leq\frac{M(x)}{\alpha\eps}.
    \label{eq:mass-volume-bound}
\end{equation}
\end{lemma}

\begin{proof}
For each \(u\in S_\eps(x)\),
\(
\abs{\sqrt{d_u}g_u(x)}\geq\alpha\eps d_u
\).
Sum over the active set.
\end{proof}

\begin{lemma}[Signed Gauss--Seidel mass decrease]
\label{lem:gs-mass-decrease}
Let \(q=D^{1/2}g(x)\), let \(u\) be processed with \(\omega=1\), and set
\(\chi=(1-\alpha)/(1+\alpha)\).  Then
\begin{equation}
    q^+=q-q_u e_u+\chi q_uAD^{-1}e_u,
    \label{eq:gs-scaled-update}
\end{equation}
and
\begin{equation}
    \norm{q^+}_1
    \leq
    \norm{q}_1-\frac{2\alpha}{1+\alpha}\abs{q_u}.
    \label{eq:gs-mass-decrease}
\end{equation}
No sign assumption on \(q\) is required.
\end{lemma}

\begin{proof}
The update formula follows by multiplying the local gradient update by
\(D^{1/2}\).  Since \(AD^{-1}e_u\geq0\) and
\(\norm{AD^{-1}e_u}_1=1\), the triangle inequality gives
\[
  \norm{q^+}_1
  \leq
  \norm{q}_1-\abs{q_u}+\chi\abs{q_u}
  =\norm{q}_1-\frac{2\alpha}{1+\alpha}\abs{q_u}.
\]
\end{proof}

\begin{theorem}[Residual-mass handoff]
\label{thm:mass-handoff}
Start the \(\omega=1\) LocSOR tail from \(x_0\).  Then
\begin{equation}
    \Work_{\rm tail}
    \leq
    \frac{1+\alpha}{2\alpha^2\eps}M(x_0).
    \label{eq:mass-tail-work}
\end{equation}
If, for a constant \(\theta>0\),
\begin{equation}
    M(x_0)\leq\theta\alpha^{3/2},
    \label{eq:mass-switch-condition}
\end{equation}
then
\begin{equation}
    \Work_{\rm tail}
    \leq
    \frac{\theta(1+\alpha)}{2\sqrt\alpha\eps},
    \qquad
    \vol(S_\eps(x_k))
    \leq\frac{\theta\sqrt\alpha}{\eps}
    \quad\text{for every tail iterate }x_k.
    \label{eq:mass-tail-accelerated}
\end{equation}
\end{theorem}

\begin{proof}
An active coordinate satisfies \(\abs{q_u}\geq\alpha\eps d_u\).  By
\cref{lem:gs-mass-decrease}, processing it decreases \(M\) by at least
\(
2\alpha^2\eps d_u/(1+\alpha)
\).  Sum this telescoping decrease to obtain
\cref{eq:mass-tail-work}.  Under \cref{eq:mass-switch-condition}, substitute
into the work bound.  The active-volume statement follows from
\cref{lem:mass-volume} and monotonicity of \(M\) during the tail.
\end{proof}

At \(x=0\), \(M(0)=\alpha\); thus
\cref{eq:mass-switch-condition} asks AESP to reduce the mass by only a factor
of order \(\sqrt\alpha\).  This is highly plausible empirically, but a
uniform degree-weighted \(\ell_1\) contraction theorem for the AESP outer
sequence is not currently available.  The objective-gap handoff is therefore
the cleaner proved route, while the residual-mass handoff is a stronger target
for future analysis and implementation.
